\section{Autonomous threshold networks admit arbitrarily many stable equilibria}
The following two theorems construct arbitrary multistability under
different restrictions on the response.
Theorem~\ref{t:siphMulti} fixes the response zero set and
allows its slope to vary, preserving every siphon and every transverse
Jacobian while creating arbitrarily many stable positive equilibria.
Indeed, for its production--loss splitting, on every siphon face,
\[
 K_\Sigma(y)
 =\partial_{x_\Sigma}(\Pi_\Phi)_\Sigma(0,y)
   (V_\Sigma)^{-1}=0,\qquad R_\Sigma(y)=0,
\]
because $\Phi'(0)=0$.
Thus unchanged, strictly subcritical boundary invasion numbers can
coexist with arbitrarily large interior multistability.
Theorem~\ref{t:KBsharp} instead fixes a slope bound and identifies a sharp
extension $\bar R_0$ of the next-generation number: below $1/L$, every response gives global exponential
convergence; above $1/L$, suitable responses give arbitrarily many stable
equilibria.
Both proofs prescribe response values at selected arguments and make
the response derivative vanish there, giving $J=-V$ at the
constructed equilibria; neither theorem requires the other.

\paragraph{Arbitrary multistability with unchanged siphons.}
Without a uniform derivative bound, polynomials $P_i$ give an exact
support criterion for constructing stable positive equilibria while
preserving every siphon and its transverse Jacobian.

\beT[Threshold networks admit arbitrarily many stable positive equilibria exactly when production reaches every coordinate]
[][siphMulti]
Fix a nonsingular $M$-matrix $V$, $b\gg0$, and polynomials
$P_i(x)=\sum_{\nu\ne0}a_{i\nu}x^\nu$ with $a_{i\nu}\ge0$,
allowing $P_i\equiv0$. Consider the production threshold model \eqref{eq:KB}, \eqref{polR}
with $h_i=0$ and a common response $\Phi$, and write $\Pi_\Phi$ for its production,
where $\mathcal C_0$ consists of the $C^\infty$ CDFs satisfying
$\Phi(q)=0$ for $q\le0$, $\Phi(q)>0$ for $q>0$, and $\Phi'\ge0$.
Every response in $\mathcal C_0$ has all derivatives zero at zero. Put
\[
 I=\{i:P_i\not\equiv0\},\qquad
 c_i=b_i\mathbf1_{\{i\in I\}},\qquad r=V^{-1}c.
\]
\begin{enumerate}
\item A positive equilibrium for any $\Phi\in\mathcal C_0$ requires $r\gg0$.
\item If $r\gg0$, every $m\ge1$ admits a common $\Phi_0\in\mathcal C_0$
 with strictly ordered positive equilibria
 \[
 x^{(k)}=t_kr,\quad 0<t_1<\cdots<t_m<1,\qquad J(x^{(k)})=-V.
 \]
 These and the origin are locally exponentially stable. A small
 common-response perturbation makes $\Phi'(q)>0$ for all $q>0$ while
 retaining $m$ strictly ordered stable positive equilibria.
\item Every allowed response has the same siphons, namely
 \[
 \Sigma\ne\varnothing:\quad V_{\Sigma\Sigma^c}=0,
 \qquad \operatorname{supp}\nu\cap\Sigma\ne\varnothing
 \quad(i\in\Sigma,\ a_{i\nu}>0).
 \]
 On each face $J_\Sigma^\perp(y)=-V_\Sigma$ for every
 admissible resident state; in particular $J(0)=-V$.
\end{enumerate}
Thus $r\gg0$ is necessary and sufficient for arbitrarily many stable positive equilibria as the common response varies
in this class; additional equilibria are not excluded.
\eeT
\begin{proof}
The common box $[0,V^{-1}b]$ is invariant since
$0\le\Pi_\Phi\le c\le b$ and $\partial_x\Pi_\Phi\ge0$.
Indeed, if $\bar x=V^{-1}b$ and $x_i=\bar x_i$, then
$f_i(x)\le-(V\bar x)_i+b_i=0$; lower-face derivatives are nonnegative.
Every equilibrium satisfies $x=V^{-1}\Pi_\Phi(x)\le r$, proving necessity.
For $r\gg0$, each $q_i(t)=P_i(tr)$, $i\in I$, increases strictly from
zero. Choose $t_m<1$ and recursively choose earlier $t_k$ so that
\[
 a_k:=\min_{i\in I}q_i(t_k)\le \widehat b_k:=\max_{i\in I}q_i(t_k)<a_{k+1}.
\]
Construct a smooth density, zero on $(-\infty,0]$ and every $[a_k,\widehat b_k]$,
with respective masses $t_1$, $t_{k+1}-t_k$, and $1-t_m$ on
$(0,a_1)$, $(\widehat b_k,a_{k+1})$, and $(\widehat b_m,\infty)$.
Positive pieces flat at their endpoints are normalized multiples of
\[
 e^{-1/((q-a)(b-q))}\quad(a<q<b),\qquad
 e^{-1/(q-\widehat b_m)-(q-\widehat b_m)}\quad(q>\widehat b_m).
\]
Their cumulative integral $\Phi_0$ belongs to $\mathcal C_0$, equals
$t_k$ on $[a_k,\widehat b_k]$, and has zero slope there. Hence
$\Pi_{\Phi_0}(t_kr)=t_kc=t_kVr$ and $J(t_kr)=-V$.
All derivatives of $\Phi_0$ vanish at zero, so the origin has the same
stable Jacobian. Perturb by
$\Phi_\varepsilon=(1-\varepsilon)\Phi_0+\varepsilon G$ with
$G(q)=e^{-1/q}$ for $q>0$ and zero otherwise. The implicit-function
theorem preserves the selected positive equilibria, strict ordering,
and stability, while $\Phi_\varepsilon'>0$ on $(0,\infty)$.

On a zero face,
$f_\Sigma(0,y)=-V_{\Sigma\Sigma^c}y+
(b_i\Phi(P_i(0,y)))_{i\in\Sigma}$ is a sum of nonnegative terms.
It vanishes identically exactly under the displayed transfer and
monomial conditions. There $P_i(0,y)=0$ for $i\in\Sigma$, and
$\Phi'(0)=0$ gives the asserted transverse matrix for every response,
including the perturbation.
\end{proof}

\paragraph{Support and scope.}
For diagonal $V$, $r\gg0$ means that every $P_i$ is nonzero; with
transfers, every coordinate must be reachable from $I$ through edges
$j\to i$ with $-v_{ij}>0$, including paths of length zero.
A common response meets any response-sharing condition within the groups of Section~\ref{s:KBred}
or overlapping siphons. Resident dynamics may change although the
siphons and transverse matrices do not.

\paragraph{Arbitrary multistability under a slope bound.}
\beT[Threshold networks realize arbitrary multistability when $\bar R_0>1/L$]
[sharp threshold $\bar R_0=1/L$][KBsharp]\label{p:KBinterpolate}\label{o:KBsharp}\label{p:KBmaxstable}\label{o:KBmaxstable}
Fix a nonsingular $M$-matrix $V$, a synaptic matrix $W\ge0$, and $L>0$.
Let $\mathcal C_L$ consist of the functions $\Phi$, smooth on $[0,\infty)$, with $\Phi\ge0$ and
$0\le \Phi'\le L$ there; the arguments $\beta Wx$ are nonnegative on the orthant. Distribution functions
with this slope bound belong to $\mathcal C_L$, and so does the identity when $L\ge1$.
Consider
\begin{equation}\label{eq:KB-sharp-family}
 \dot x=-Vx+\Phi(\beta Wx),\qquad x\in\mathbb R_+^{n_S},
 \quad \Phi\in\mathcal C_L,\quad\beta\ge0,
\end{equation}
with the same scalar response applied to every coordinate. Put
\[
 r_0=\rho(WV^{-1})=\rho(V^{-1}W),\qquad
 \bar R_0=\rho(\beta WV^{-1})=\beta r_0.
\]

\begin{enumerate}
\item If $\bar R_0<1/L$, every $\Phi\in\mathcal C_L$ gives a unique
globally exponentially attracting equilibrium.
\item If $\bar R_0>1/L$, for every $m\ge2$ there exists
$\Phi\in\mathcal C_L$ giving at least $m$ distinct locally exponentially
stable equilibria. One can choose $\Phi(q)=0$ for $q\le0$ and obtain
\[
 x^{(k)}=t_kr,\quad 0=t_0<t_1<\cdots<t_{m-1},
 \qquad r\ge0,\quad r\ne0,\qquad J(x^{(k)})=-V.
\]
Alternatively, one can choose $0<\Phi(q)<1$ and $\Phi'(q)>0$ on
$\mathbb R$ and obtain $m$ strictly positive stable equilibria.
The response may depend on $m$, and in both cases it is a distribution function, so
the construction stays within bounded responses.
\end{enumerate}
Thus $\bar R_0=1/L$ is the sharp threshold for coexistence of equilibria within $\mathcal C_L$.
No irreducibility is required; uniqueness at equality is not asserted.
A common additive stimulus can be absorbed into $\Phi$, provided $\Phi$ is defined and nonnegative on the
arguments attained.

Thus the stable count is unbounded over this response class, even
for $n_S=1$. With $\Phi>0$ and $\Phi'>0$, the interaction graph remains
fixed and there are no nonempty siphons. In the zero-onset construction, the selected equilibria are ordered
componentwise and lie on the ray $\{tr:t\ge0\}$.
This ray need not be invariant. The full-support perturbation
preserves their number and stability, but need not preserve
their location on that ray.
\eeT

\beR[$\bar R_0$ extends the next-generation number][]
[NGMg]
Write
\[
 J(x)=F(x)-V,\qquad
 F(x)=\beta\operatorname{Diag}(\Phi'(\beta Wx))W.
\]
Since $0\le F(x)\le\beta LW$,
\[
 R(x):=\rho(F(x)V^{-1})\le L\bar R_0.
\]
At an equilibrium $x_*$, the NGM matrix criterion gives
$J(x_*)$ Hurwitz if and only if $R(x_*)<1$. If $\Phi(0)=0$, the origin is an equilibrium and
$R(0)=\rho(\beta\Phi'(0)WV^{-1})=\Phi'(0)\bar R_0$ is the classical $R_0$, so the classical test is $\bar R_0<1/\Phi'(0)$, which is $\bar R_0<1/L$ when $L=\Phi'(0)$.
In general the test $\bar R_0<1/L$ is uniform over all states and
responses in $\mathcal C_L$, and the bar marks that $\bar R_0$ is not the reproduction number
of a specified equilibrium. The condition $\bar R_0<1/L$ ensures global
exponential convergence; $\bar R_0>1/L$ permits arbitrarily many stable
equilibria through a suitable choice of response.
\eeR

\begin{proof}
\emph{Interpolation used in the construction.}
Let $0<c<L$ and $0<q_1<\cdots<q_d<1/c$. There is a $C^\infty$ CDF
$\Phi_0:\mathbb R\to[0,1]$ satisfying
\[
 \Phi_0(q)=0\ (q\le0),\qquad 0\le \Phi_0'(q)<L,
 \qquad \Phi_0(q_j)=cq_j,\quad \Phi_0'(q_j)=0.
\]
Its density has compact support in $(0,\infty)$ and vanishes near each
$q_j$, including $\Phi_0'(0)=0$.


Choose $\eta\in C_c^\infty(0,1)$, $\eta\ge0$, with $\int\eta=1$ and
$\|\eta\|_\infty<L/c$. Such densities have supremum arbitrarily close
to one. Put $q_0=0$ and choose $a>(1-cq_d)\|\eta\|_\infty/L$. Set
\[
 \varphi(q)=\begin{cases}
 c\eta((q-q_{j-1})/(q_j-q_{j-1})),&q_{j-1}<q<q_j,\\
 (1-cq_d)\eta((q-q_d)/a)/a,&q_d<q<q_d+a,\\
 0,&\text{otherwise}.
 \end{cases}
\]
The pieces have masses $c(q_j-q_{j-1})$ and $1-cq_d$, vanish near
their endpoints, and have supremum less than $L$. Their cumulative
integral is the required response.

\emph{Contraction below the bound.}

For $\bar R_0<1/L$, the matrix $V-\beta LW$ is a nonsingular $M$-matrix, because
$\rho(\beta LWV^{-1})<1$. Put
\[ p=(V-\beta LW)^{-1}\mathbf1\gg0 .\]
The bound $\partial_x\Pi\le\beta LW$ gives $J(x)\le-V+\beta LW$ entrywise, so
$J(x)p\le-\mathbf1$ for every $x$. Since $J(x)$ is Metzler, this gives contraction
in the weighted maximum norm $\max_i|x_i|/p_i$, at rate $1/\max_ip_i$. The vector field is globally
Lipschitz on the orthant, whose faces $x_i=0$ point inward, so nonnegative solutions are global. The
time-$t$ maps are strict contractions of the complete orthant, so they have one common fixed point,
which is the unique equilibrium, and every solution converges to it exponentially.
\emph{Stable equilibria above the bound.}
For $\bar R_0>1/L$, choose a nonnegative Perron vector
\[
 V^{-1}Wr=r_0r,\qquad r\ne0,\qquad Wr=r_0Vr\ne0,
\]
and put $c=1/\bar R_0<L$. Choose the positive $t_k$ so that
$\max_i\beta t_{m-1}(Wr)_i<1/c$. Apply the interpolation construction above to all distinct positive arguments
$\beta t_k(Wr)_i$. At these and at zero,
\[
 \Phi_0(\beta Wt_kr)=c\beta t_kWr=t_kVr,
 \qquad
 J(t_kr)=-V+\beta\operatorname{Diag}(\Phi_0'(\beta Wt_kr))W=-V.
\]
The $m$ equilibria are exponentially stable and lie in
$[0,V^{-1}\mathbf1]$.

\emph{Strictly positive equilibria and full support.}
Set
\[
 G(q)=(1+e^{-Lq})^{-1},\qquad
 \Phi_\varepsilon=(1-\varepsilon)\Phi_0+\varepsilon G.
\]
Since $0<G'\le L/4$, this remains in $\mathcal C_L$ and is strictly
increasing. The implicit-function theorem continues the equilibria, with
\[
 \left.\frac{d x^{(k)}(\varepsilon)}{d\varepsilon}\right|_0
 =V^{-1}G(\beta Wx^{(k)})-x^{(k)}.
\]
The first term is strictly positive: $V^{-1}\ge0$ has no zero row.
Every initially zero coordinate therefore becomes positive, while
positive coordinates remain positive. Disjoint neighbourhoods and
hyperbolicity preserve the count and stability for one small
$\varepsilon>0$. Replacing $\Phi(q)$ by $\Phi(q+h)$ absorbs a common stimulus
$h\mathbf1$. Full support excludes zero faces; for these constructions $\beta>0$,
and the off-diagonal Jacobian entries are positive exactly when $-v_{ij}>0$ or $w_{ij}>0$.
\end{proof}

\beR[Fixed production coefficients and the scope of the stable-equilibrium count]
[fixed coefficients and counting scope][KBfixedamp]\label{r:KBcountscope}
For $-Vx+\operatorname{Diag}(b_i)\Phi(\beta Wx)$, $b_i>0$, replace $r_0$
by $\rho(\operatorname{Diag}(b_i)WV^{-1})$ and use a Perron vector of
$V^{-1}\operatorname{Diag}(b_i)W$. The interpolation values remain
$\Phi(q)=q/(\beta r_0)$, and the selected Jacobians remain $-V$.


The response in Theorem~\ref{t:KBsharp} varies with the requested count. Finite bounds require
additional restrictions, such as a fixed distribution family, density
oscillation bounds, or fixed polynomial degree. This construction
preserves neither unimodality nor fixed-degree SIS structure, and does
not assert infinitely many hyperbolic equilibria for one response.
\eeR

\beC[Subcritical threshold networks entrain to periodic stimuli][periodic entrainment][Ent]
Let $V$ be a nonsingular $M$-matrix, $W\ge0$, $\beta\ge0$, and let $\Phi:\mathbb R\to[0,1]$ be $C^1$
with $0\le \Phi'\le L$, where $L>0$. Apply $\Phi$ coordinatewise and let
$h:\mathbb R\to\mathbb R^{n_S}$ be continuous and $T$-periodic, with $T>0$. If
\[
 \bar R_0=\beta\rho(V^{-1}W)<\frac1L,
\]
then
\[
 \dot x=-Vx+\Phi(\beta Wx+h(t))
\]
has a unique $T$-periodic solution $x_{\rm per}(t)$ in the nonnegative orthant. Every nonnegative
solution satisfies
\[
 \|x(t)-x_{\rm per}(t)\|_{\infty,p}
 \le e^{-ct}\|x(0)-x_{\rm per}(0)\|_{\infty,p},
\]
where $p=(V-\beta LW)^{-1}\mathbf1\gg0$, $c=\min_i p_i^{-1}>0$ and
$\|z\|_{\infty,p}=\max_i|z_i|/p_i$. The forcing need not be small or slowly varying.
This is the standard entrainment consequence of contraction
\cite{RussoDiBernardoSontag2010,Bullo2026}.
\eeC
\begin{proof}
As in the proof of Theorem~\ref{t:KBsharp}, the Jacobian is Metzler and obeys
\[
 J(x,t)=-V+\beta\operatorname{Diag}(\Phi'(\beta Wx+h(t)))W\le -V+\beta LW,
 \qquad J(x,t)p\le-\mathbf1\le-cp.
\]
The weighted maximum-norm estimate therefore holds for any pair of solutions with the same stimulus.
Put $q=V^{-1}\mathbf1\gg0$. On $x_i=0$, $x\ge0$, the vector field points into the orthant. On the
upper $i$th face of $[0,q]$,
\[
 f_i(x,t)\le-(Vq)_i+1=0,
\]
because $-V$ is Metzler. Thus $[0,q]$ is invariant, and its time-$T$ map has Lipschitz constant at
most $e^{-cT}<1$. Banach's theorem gives a periodic solution. Bounded production ensures global
existence, and contraction proves attraction from every nonnegative initial state and uniqueness
among all nonnegative $T$-periodic solutions. The minimal period may divide $T$.
\end{proof}

\paragraph{Arbitrary multistability with convergence of every trajectory.}
Theorem~2 of \cite{OfirOvseevichMargaliot2024} complements Theorem~\ref{t:KBsharp} by providing
convergence without uniqueness. For
\[
 \dot x=-dx+\Phi(\beta Wx),\qquad d>0,\qquad \Phi\ \text{bounded},\quad 0\le \Phi'\le L,
\]
their sufficient condition for $2$-contraction is
\[
 \beta^2L^2\bigl(\sigma_1(W)^2+\sigma_2(W)^2\bigr)<2d^2,
\]
where $\sigma_1(W)\ge\sigma_2(W)$ are the two largest singular values. Since $\Phi$ is bounded,
every trajectory then converges to an equilibrium, which may depend on the initial state.

This condition is compatible with $\bar R_0>1/L$. Let $W=I+C$, with $C$ the cyclic permutation matrix
of three coordinates, so that $\rho(W)=\sigma_1(W)=2$ and $\sigma_2(W)=1$. If
\[
 \frac12<\frac{\beta L}d<\sqrt{\frac25},
\]
then, for every $m\ge2$, Theorem~\ref{t:KBsharp} provides a common smooth distribution function
with $0<\Phi<1$ and $0<\Phi'\le L$ giving at least $m$ strictly positive locally exponentially stable
equilibria, and every nonnegative trajectory converges to an equilibrium. For rank-one
$W=w\,ee^\top$ with $e\gg0$, convergence needs no such criterion: $z=e^\top x$ satisfies a scalar
autonomous equation, and each coordinate is then driven by a convergent production term.

% Preserved original lifting theorem; merged into commonRed.
% \beT[Folds and cusps of invariant reductions persist in the full network]
% [lifting reduction bifurcations][KBbiflift]
% Let $f$ be a $C^k$ family, $k\ge3$, and let
% $H\in\mathbb R^{N\times m}$ be the indicator matrix of a partition
% $V_1,\ldots,V_m$ into nonempty groups:
% $H_{ia}=1$ for $i\in V_a$ and $H_{ia}=0$ otherwise. Suppose
% \[
%  f(Hy;\theta)=H\widehat f(y;\theta)
% \]
% for all nearby $y,\theta$. At an equilibrium $x_*=Hy_*$, put
% $J=\partial_xf(x_*;\theta_*)$ and
% $\widehat J=\partial_y\widehat f(y_*;\theta_*)$.
% Assume either:
% \begin{enumerate}
% \item $J$ is irreducible Metzler and $s(\widehat J)=0$; or
% \item $\widehat J$ has a zero eigenvalue and $J$ has a simple zero
%  eigenvalue with all other eigenvalues in the open left half-plane.
% \end{enumerate}
% In the second alternative, $H$ may more generally be any fixed matrix
% of full column rank.
% Then, with the parameters included as constant variables, a reduced
% center manifold lifts to a full center manifold.
% A nondegenerate reduced saddle-node or generically unfolded cusp is
% therefore the same full-network bifurcation, at the same parameter
% location; stability of its center branches determines their full stability.
% \eeT
% \begin{proof}
% The identity makes $\operatorname{im}H$ invariant for every nearby
% parameter and gives $JH=H\widehat J$. In the first case,
% Theorem~\ref{t:Spe} yields $s(J)=s(\widehat J)=0$; irreducible
% Metzler theory gives a simple zero and a strictly stable complement.
% The second case assumes these spectral properties directly.
% Because $H$ is injective, the reduced zero eigenvalue is simple,
% its kernel maps onto $\ker J$, and all remaining reduced eigenvalues
% have negative real parts.
%
% Use the embedding identity of Proposition~\ref{p:LinRed}.
% Choose any linear complement to $\operatorname{im}H$ and write
% $x=Hy+Zz$, with $z=0$ the synchrony subspace. In these coordinates,
% including $\dot\theta=0$, the linearization has the form
% \[
%  \begin{pmatrix}
%  \widehat J&\partial_\theta\widehat f&*\\
%  0&0&0\\
%  0&0&A_\perp
%  \end{pmatrix}
%  \quad\text{in the order }(y,\theta,z).
% \]
% The lower-left entries vanish because $z=0$ is invariant for every
% nearby $\theta$. The matrix $A_\perp$ has strictly negative spectral
% abscissa, so the full generalized center subspace lies in $z=0$,
% including the parameter directions and their possible generalized
% eigenvectors. A center manifold of the extended reduced system thus
% lifts under $(y,\theta)\mapsto(Hy,\theta)$ to an invariant manifold
% tangent to the full center subspace. Its dynamics are unchanged.
% The complementary eigenvalues remain strictly stable nearby, proving
% both the bifurcation and stability assertions.
% \end{proof}
% 